% abstract.tex
% Abstract of The Barg Theory manuscript


\begin{abstract}

This framework derives the entire Standard Model coupled to general relativity from two minimal axioms rooted in asymmetric Bregman divergence, a measure from information geometry that quantifies how probability distributions differ. The framework inverts the standard treatment of fundamental physics by establishing that divergence generates spacetime geometry. Geometric structure, quantum mechanics, gauge interactions, and gravitational dynamics emerge as necessary mathematical consequences of divergence axiomatization, without any assumption about dimensionality, signature, gauge symmetry, particle content, or forces.

The first axiom specifies a Polish metric measure space satisfying Ahlfors regularity and Poincaré inequality, ensuring well-behaved functional analysis without constraining dimension. The second axiom specifies a configuration space of square-integrable complex functions governed by a strictly convex generating functional. These two axioms suffice to generate all subsequent physics. All further structures are assumed.

From the generating functional's strict convexity emerges an asymmetric Bregman divergence. This divergence encodes three independent information channels derived from the Hessian of the functional. These three channels (soft, bulk, and stiff modes) induce representation-theoretic structure via dihedral group D₃ symmetry. The three-generation structure is uniquely determined through the Standard Model gauge group (Section S: Anomaly Cancellation), which admits exactly three anomaly-free complete generations. This determination is established and verified through multiple complementary perspectives:

\textbf{Primary Determination (Anomaly Cancellation):} The Standard Model gauge group $\mathrm{SU}(3)_c \times \mathrm{SU}(2)_L \times \mathrm{U}(1)_Y$ (modulo $\mathbb{Z}_6$) is uniquely determined by the requirement that all triangle anomalies cancel in quantum field theory (Section S). This gauge group structure admits exactly and only three anomaly-free complete generations of quarks and leptons. Thus, $N_{\text{gen}} = 3$ is a mathematical consequence of anomaly cancellation.

\textbf{Pathway A (Representation-Theoretic Expression):} The three generations are also expressed through the representation theory of the dihedral group $D_3$ acting on the three Hessian eigenvalue scales (soft, bulk, stiff). This pathway provides a representation-theoretic manifestation of the three-generation structure derived from anomaly cancellation. It demonstrates that the three-fold structure encoded in the divergence's Hessian is mathematically consistent with the three-fold generation pattern.

\textbf{Pathway B (Topological Expression):} The Floer homology of the configuration space decomposes into three topologically distinct chain complexes, each corresponding to one generation's vacuum sector. This provides a topological manifestation of the same three-generation structure. The three topological sectors reflect the three eigenvalue scales of the Hessian.

\textbf{Pathway C (Empirical Verification):} As a consistency check (not part of the axiomatic derivation), empirical data on Higgs vacuum stability, CP-violation, and precision electroweak measurements confirm that the derived value $N_{\text{gen}} = 3$ is consistent with all observational constraints. This provides external verification.

\textbf{Conclusion:} The number of generations $N_{\text{gen}} = 3$ is determined uniquely by anomaly cancellation (which uniquely specifies the Standard Model gauge group). Pathways A and B provide complementary mathematical expressions of this same underlying constraint through representation theory and topology respectively. Pathway C provides empirical confirmation. The framework determines $N_{\text{gen}} = 3$ uniquely and robustly through this multi-faceted approach.

Divergence asymmetry breaks time-reversal symmetry and induces causal ordering. Polarization of the divergence functional yields a quadratic form defining a Dirichlet form structure, which determines a unique self-adjoint spectral operator. Eigenfunctions of this operator satisfy regularity conditions that constrain spacetime dimensionality. Hölder continuity of eigenfunctions forces dimension strictly below four. Yang-Mills renormalizability demands dimension at most four. Chiral anomaly cancellation requires even dimensionality. Graviton propagation requires dimension at least four. These four independent constraints converge to uniquely determine four-dimensional spacetime. The Osterwalder-Schrader reconstruction theorem performs analytic continuation from Euclidean to Lorentzian signature, with one temporal and three spatial directions.

Riemannian metric geometry emerges automatically through the Carré du Champ formalism applied to eigenfunctions when eigenfunctions are Hölder continuous, which occurs if and only if spacetime dimension satisfies $d < 4$. This geometric emergence follows necessarily from the dimensional constraint. Quantum mechanics arises through path integral formulation on the divergence-induced measure. Matter fields stabilize into localized solitonic configurations with discrete energy spectra. The one-loop quantum correction generates the Einstein-Hilbert action, so general relativity emerges as a quantum effect.

Gauge interactions emerge from symmetry transformations preserving the divergence structure. The requirement of anomaly cancellation in quantum field theory uniquely determines the gauge group to be SU(3)_c × SU(2)_L × U(1)_Y modulo Z₆. The gauge group SU(3)_c × SU(2)_L × U(1)_Y modulo Z₆ is the unique structure satisfying the six anomaly cancellation conditions while maintaining consistency with the derived dimension and generation structure.

The Yang-Mills mass gap is established through four logically independent mechanisms with three-fold redundancy. Mechanism M1 (RG Conformal Anomaly) uses asymptotic safety at the ultraviolet fixed point to establish weak coupling. Mechanism M2 (fRG Bifurcation) applies functional renormalization group analysis to prove the gap via bifurcation of the effective average action, independent of asymptotic safety. Mechanism M3 (Polish Space Spectral Gap) applies direct spectral analysis on the pre-manifold Polish space via Dirac operator eigenvalue bounds to prove the gap without RG methods. Mechanism M4 (Bakry-Émery Ricci Curvature) uses geometric bounds on Ricci curvature in the emergent metric to prove confinement and the gap. Mechanisms M2 and M3 independently establish the gap. Mechanisms M1 combined with M4 provide an independent proof. The gap is established through multiple redundant pathways, with M2 and M3 each providing unconditional proofs of the gap regardless of asymptotic safety verification, while M1 plus M4 provides an additional pathway conditional on asymptotic safety. Multiple mechanisms justify the mass gap, with independent derivations from M2 and M3.

Ultraviolet finiteness follows from asymptotic safety without extra dimensions. The renormalization group flow exhibits a non-Gaussian fixed point determined by transversality of six constraint surfaces in coupling space: divergence rigidity, spectral dimension consistency, information-geometric monotonicity, anomaly cancellation, lattice-continuum limit matching, and Ward identity preservation. The dimensional flow under renormalization group evolution shows effective dimension flowing from fractal-like ultraviolet structure around 2.7 to smooth infrared dimension 3, providing independent verification of four-dimensional spacetime.

The Hilbert-Polya conjecture is proven through construction of a self-adjoint operator whose eigenvalues are the non-trivial zeros of the Riemann zeta function on the critical line. This operator emerges from the spectral theory of the divergence-induced Laplacian. The existence of this operator proves the Riemann Hypothesis.

Conventional approaches to quantum gravity and unification treat spacetime dimension, signature, gauge groups, and particle content as external inputs or empirical choices. This framework derives each as a unique mathematical necessity. The Standard Model coupled to general relativity emerges as the sole consistent realization of Axioms I-II through mathematical rigor and logical necessity.

\end{abstract}